%\makeglossaries
\makenoidxglossaries

% Danh mục thuật ngữ và từ viết tắt
\newglossaryentry{HCI}{
    type=\acronymtype,
    name={HCI},
    description={Tương tác người - máy (Human - Computer Interaction)},
    first={Tương tác người - máy (Human - Computer Interaction - HCI)}
}
\newglossaryentry{NUI}{
    type=\acronymtype,
    name={NUI},
    description={Giao diện người dùng tự nhiên (Natural User Interface)},
    first={Giao diện người dùng tự nhiên (Natural User Interface - NUI)}
}
\newglossaryentry{UDP}{
    type=\acronymtype,
    name={UDP},
    description={Giao thức gói người dùng (User Datagram Protocol)},
    first={Giao thức gói người dùng (User Datagram Protocol - UDP)}
}
\newglossaryentry{TFLite}{
    type=\acronymtype,
    name={TFLite},
    description={Phiên bản tối ưu hóa TensorFlow cho thiết bị di động và nhúng (TensorFlow Lite)},
    first={TensorFlow Lite (TFLite)}
}
\newglossaryentry{MLP}{
    type=\acronymtype,
    name={MLP},
    description={Mạng nơ-ron truyền thẳng nhiều lớp (Multi-Layer Perceptron)},
    first={Mạng nơ-ron truyền thẳng nhiều lớp (Multi-Layer Perceptron - MLP)}
}
\newglossaryentry{LSTM}{
    type=\acronymtype,
    name={LSTM},
    description={Mạng bộ nhớ ngắn hạn dài (Long Short-Term Memory)},
    first={Mạng bộ nhớ ngắn hạn dài (Long Short-Term Memory - LSTM)}
}
\newglossaryentry{FPS}{
    type=\acronymtype,
    name={FPS},
    description={Số khung hình hiển thị trên mỗi giây (Frames Per Second)},
    first={Số khung hình trên giây (Frames Per Second - FPS)}
}
\newglossaryentry{API}{
    type=\acronymtype,
    name={API},
    description={Giao diện lập trình ứng dụng (Application Programming Interface)},
    first={Giao diện lập trình ứng dụng (Application Programming Interface - API)}
}
\newglossaryentry{CSV}{
    type=\acronymtype,
    name={CSV},
    description={Tệp giá trị phân tách bằng dấu phẩy (Comma-Separated Values)},
    first={Tệp giá trị phân tách bằng dấu phẩy (Comma-Separated Values - CSV)}
}